\section{Schur pairings before electric dressing}\label{sec:schur}
\subsection{The conformal electric benchmark}
Consider the four-dimensional $\mathcal N=2$ $SU(2)$ theory with four
hypermultiplets, with $0<q<1$ and trivial flavour fugacities. Let
\begin{equation}\label{eq:pochhammer}
 (a;Q)_N=\prod_{j=0}^{N-1}(1-aQ^j),\qquad
 (a;Q)_\infty=\prod_{j\geq0}(1-aQ^j).
\end{equation}
The electric-sector spherical norm from \cite[Section 5.1.5]{Schur}
specializes to the finite positive measure
\begin{align}
 \dd\nu_q(\theta)&=Z_q^{-1}\omega_q(\theta)\frac{\dd\theta}{2\pi},
 \qquad Z_q=\int_0^{2\pi}\omega_q(\theta)\frac{\dd\theta}{2\pi},
 \notag\\
 \omega_q(\theta)&=4\sin^2\theta\,
 \frac{|(q^2e^{2i\theta};q^2)_\infty|^4}
      {|(-qe^{i\theta};q^2)_\infty|^{16}}.
 \label{eq:su2-measure}
\end{align}
The Wilson operator acts as $W=2\cos\theta$. Normalizing by $Z_q$
sets the vacuum norm to one; the overall physical normalization would
also have to be fixed by a full matching theorem.

For the natural source $a_f(W)=\int f(x)e^{ixW}\dd x$, the entire
pairing is
\begin{equation}\label{eq:wilson-norm}
 B_q(f,g)=\int_0^{2\pi}
 \overline{\widehat f(-2\cos\theta)}\widehat g(-2\cos\theta)
 \dd\nu_q(\theta).
\end{equation}
Its kernel $\int e^{2ir\cos\theta}\dd\nu_q(\theta)$ is entire in
$r$, because the spectral variable lies in $[-2,2]$. It cannot produce
separated prime atoms, and its bounded input norm contradicts
Proposition~\ref{prop:bounded-source}. This conclusion concerns this
electric preparation only.

\subsection{Local difference factors and the full spherical vector}
The elementary factor
\begin{equation}\label{eq:chi}
 \chi_q(v)=(-qv;q^2)_\infty^{-1}
 \quad\hbox{satisfies}\quad
 \frac{\chi_q(q^2v)}{\chi_q(v)}=1+qv.
\end{equation}
It occurs in the elementary $q$-Weyl example of
Ambrosino--Gaiotto \cite[Section 3.1.1]{RG}. The ratio follows by
cancelling the products in \eqref{eq:pochhammer}.
For a trial substitution $q=p^{-1/2}$ and $v=-e^{i\tau\log p}$, let
$J_p(\tau)=1-p^{-1/2}e^{i\tau\log p}$. Then
\begin{equation}\label{eq:prime-phase}
 2\im\partial_\tau\log J_p(\tau)
 =-2(\log p)\sum_{m\geq1}p^{-m/2}\cos(m\tau\log p).
\end{equation}
The series converges absolutely. The unshifted logarithm of $\chi_q$
instead has coefficients with a factor $(1-q^{2m})^{-1}$. Thus the
ratio in \eqref{eq:chi}, rather than the factor itself, supplies the
geometric coefficient.

This calculation has two immediate limitations. A phase derivative
is not a state norm. In addition, varying $q$ with $p$ varies the
quantization parameter, and is not a construction in one fixed theory.
The dressed source in the next section addresses both points for the
positive local reference. In the full massless four-flavour benchmark,
the spherical vector has instead the complete ratio
\begin{equation}\label{eq:full-ratio}
 \frac{\varphi_0(q^2v)}{\varphi_0(v)}
 =\frac{(1+qv)^8}
 {(1-v^2)(1-q^2v^2)^2(1-q^4v^2)}.
\end{equation}
All matter and vector factors are retained here. Extracting one
numerator factor would change the model calculation. Complex holonomy
shifts and changes of magnetic charge also require an explicit source
map before they can represent translation in the real arithmetic
coordinate.

\subsection{The common-state RG formulation and the undressed test}
The RG formulation \cite[Sections 1--2]{RG} transports a UV line $a$
to a framed expansion $F_a$ acting on a common spectrum-generator state
$\mathscr S$. For charge lattice $\Gamma$, the relevant quantum torus
and pairing take the form
\begin{equation}\label{eq:rg-pairing}
 X_\gamma X_\eta=q^{\langle\gamma,\eta\rangle}X_{\gamma+\eta},
 \qquad I(a,b)=c_q\ip{F_a\mathscr S}{F_b\mathscr S}_{\ell^2(\Gamma)},
 \quad c_q=(q^2;q^2)_\infty^{2r_{\rm g}}>0.
\end{equation}
Here $\langle\gamma,\eta\rangle$ is the integral charge pairing and
$r_{\rm g}$ is the gauge rank. The same state enters both the
dualization rule and the norm. General domain and construction claims
in this formalism include conjectural input; only the explicitly
defined elementary Hilbert-space calculations below are used as
unconditional statements. An isometric RG description preserves the
norm of a given source and cannot itself repair a mismatch.

On one charge ray, take the bilateral shift $Xe_n=e_{n+1}$ on
$\ell^2(\Z)$ and the vector
\begin{equation}\label{eq:bare-state}
 \mathscr S=\sum_{n\geq0}c_ne_n,\qquad
 c_n=\frac{(-q)^n}{(q^2;q^2)_n},\qquad
 Z=\sum_{n\geq0}|c_n|^2<\infty.
\end{equation}
These are the Fourier coefficients of $\chi_q$. The recurrence
$(1-q^{2n})c_n=-qc_{n-1}$ gives the complete norm identity
\begin{equation}\label{eq:bare-norm}
 \norm{(1+qX)\mathscr S}^2
 =(1+q^2)Z+2q\re\ip{\mathscr S}{X\mathscr S}
 =\sum_{n\geq0}q^{4n}|c_n|^2.
\end{equation}
Indeed each coefficient of $(1+qX)\mathscr S$ is $q^{2n}c_n$,
including the coefficient at $n=0$.

The actual normalized overlaps do not give $q^m$. For $m\geq1$,
\begin{equation}\label{eq:bare-overlap}
 (-1)^m\frac{\ip{\mathscr S}{X^m\mathscr S}}Z
 =\frac{q^m}Z\sum_{n\geq0}
 \frac{|c_n|^2}{\prod_{j=1}^m(1-q^{2(n+j)})}>q^m.
\end{equation}
Every summand is positive and each denominator is strictly less than
one. Convergence follows from $(q^2;q^2)_n\to(q^2;q^2)_\infty>0$.
Thus a correct local difference equation does not by itself determine
the desired physical return coefficient. This exact failure motivates
changing the prepared state while retaining the positive Hilbert metric.
